\section{Related Work}\label{sec:related}

\paragraph{Shared passed sets in explicit-state model checking.}
Stern and Dill~\cite{Stern01} partitioned the state space across processors and
routed successors to their owner; multi-core \textsc{Spin}~\cite{Holzmann07} and
\textsc{DiVinE}~\cite{Barnat12} refined that shape for shared memory, and
\ltsmin{}~\cite{Laarman10,Kant15} replaced it with a single lockless hash table
that every worker inserts into, removing the routing phase.  Our striped passed
set follows that choice of a shared structure over a static partition.  What does
not carry over is the operation: a lockless hash table implements
insert-if-absent, which is idempotent and pure in its effect, whereas
\cref{alg:update} retires and replaces stored zones.  This is why the engine uses
a lock per stripe and a shared/exclusive split rather than a compare-and-swap on
a slot, and why \cref{thm:probe} has to be proved.

\paragraph{Parallel and distributed reachability for timed automata.}
Behrmann et al.~\cite{Behrmann00,Behrmann05} distributed \uppaal{}'s search and
observed that the search order changes the number of explored states, an effect
that reappears here as the reason for keeping strict level synchronisation.
Dalsgaard et al.~\cite{Dalsgaard12} combine the \opaal{} front end with
\ltsmin{}'s shared table to obtain a multi-core timed-automaton reachability
tool, and identify the interaction between subsumption and a shared table as the
central difficulty.  Our contribution to that point is \cref{lem:monotone} and
its consequences: what makes a shared subsumption test usable without validation
is that coverage is monotone, and once that is available the read-mostly nature
of the operation can be exploited directly.

\paragraph{Subsumption and antichains.}
Storing sets of incomparable elements under a subsumption preorder rather than
individual elements is the antichain technique of De Wulf et al.~\cite{Wulf06}.
For timed automata the relevant preorders are the $\preceq_{LU}$-based
abstractions of Behrmann et al.~\cite{Behrmann06} and Herbreteau et
al.~\cite{Herbreteau12}, implemented by \uppaal{} and
\textsc{TChecker}~\cite{Herbreteau13}; extrapolation~\cite{Daws98,Behrmann06} is
the older route to finiteness, and the engine chooses between the two per run
rather than per tool.  Termination arguments for antichain-based searches
normally rest on the stored set being an antichain, which convex merging breaks;
\cref{thm:term} replaces that argument with one that bounds insertions by
abstraction classes using only monotonicity of coverage.

\paragraph{Compact zone representations.}
Storing zones as minimal constraint systems is due to Larsen et
al.~\cite{LarsenLPY97}.  The representation is chosen for bytes resident, and
the cost it moves is the cost of a query, which must rebuild a dense matrix.
\Cref{lem:minclosure} avoids the rebuild for the query that dominates by deciding
inclusion against the generators, and the dense mirror buys the dense form back
only for the stored sets whose measured query traffic pays for it.

\paragraph{Determinism.}
Blelloch et al.~\cite{Blelloch12} argue for internally deterministic parallel
algorithms, whose intermediate state does not depend on the schedule.  This
engine is not one, and cannot cheaply be made one, because convex merging is not
confluent; \cref{thm:verdict} isolates what is deterministic --- the verdict and
the set of covered symbolic states --- and every mechanism is checked against
that property rather than against reproducibility of the passed set.
